\documentclass[conference]{IEEEtran}

\usepackage{cite}
\usepackage{url}
\usepackage{graphicx}
\usepackage{amsmath}
\usepackage{listings}
\usepackage{xcolor}
\usepackage{float}
\usepackage{hyperref}

\title{
Investigating and Implementing LibreRouter and LibreMesh:
Virtualized and Raspberry Pi-Based Mesh Networking
}

\author{
\IEEEauthorblockN{\small
Bastas, Dimitrios (SUPREMER LEFTIEUENT COMMANDER)\\
Kamilos, Dimitrios (SUPREMEISH RIGHTIEUNENT COMMANDER)\\
Robert Leyba (LT)\\
Rachel Goldstein (LT)\\
Margaret Moore (LTJG)
}
\IEEEauthorblockA{
EC3730 Final Project
}
}

\begin{document}

\maketitle

\begin{abstract}

This project investigates LibreRouter and LibreMesh, open-source community networking technologies built on top of OpenWrt. The project explores the architecture, implementation, vulnerabilities, and deployment challenges associated with decentralized mesh networking systems.

A virtualized OpenWrt environment was developed using VMware Workstation Pro to simulate a four-node routed mesh backbone topology. The environment established reproducible experimentation capabilities for distributed routing analysis, virtualization behavior, and future LibreMesh deployment.

The project will later extend the deployment onto Raspberry Pi hardware platforms for physical mesh networking experimentation and wireless performance evaluation.

\end{abstract}

\section{Introduction}

Community mesh networks provide decentralized communication infrastructure capable of operating independently of centralized Internet Service Providers (ISPs). LibreRouter and LibreMesh were developed to support resilient, community-owned wireless infrastructure using open-source software and commodity hardware.

This project investigates the deployment and operation of LibreMesh-based networking systems using both virtualized and physical implementations. The primary goals include evaluating routing behavior, deployment complexity, security considerations, network resilience, and practical limitations of community mesh networking technologies.

In addition to traditional research analysis, this report maintains a reproducible engineering notebook style documenting implementation steps, troubleshooting observations, operational anomalies, and experimental outcomes throughout the project lifecycle.

\section{Background Research}

\subsection{OpenWrt}

OpenWrt is an open-source Linux-based operating system designed primarily for embedded networking devices such as routers and access points. Unlike traditional vendor firmware, OpenWrt provides a writable filesystem, package management system, and modular architecture allowing users to customize routing, firewall, wireless, and network management functionality.

The OpenWrt ecosystem supports x86 virtualized environments in addition to physical embedded hardware, making it suitable for experimental mesh networking research.

\subsection{LibreMesh}

LibreMesh is a framework built on top of OpenWrt that simplifies deployment of decentralized mesh networking infrastructure. The platform automates many aspects of wireless mesh configuration while supporting routing protocols such as BATMAN-adv and BMX7.

LibreMesh was designed to support community-operated networks where nodes dynamically discover neighboring devices and establish resilient communication paths without centralized infrastructure.

\subsection{LibreRouter}

LibreRouter is an open-hardware networking platform developed specifically for community mesh deployments. The platform combines OpenWrt-based firmware with wireless hardware optimized for long-range and community-managed networking applications.

\section{Literature Review and Security Analysis}

\subsection{Identified Weaknesses and Vulnerabilities}

Several security concerns and operational limitations were identified during preliminary literature review:

\begin{itemize}
    \item Firmware update vulnerabilities in OpenWrt package management systems
    \item Mesh routing attacks including route poisoning and malicious node participation
    \item Wireless sniffing and unauthorized network access
    \item Resource exhaustion and denial-of-service attacks
    \item Complexity of maintaining decentralized routing infrastructure
\end{itemize}

\section{Virtualized OpenWrt Deployment}

\subsection{Virtualization Environment}

VMware Workstation Pro 17 was selected as the virtualization platform for the initial implementation phase. Virtualization provided rapid deployment, rollback capability through snapshots, reproducible experimentation, and simplified troubleshooting prior to deployment on physical Raspberry Pi hardware.

\subsection{OpenWrt Image Preparation}

The OpenWrt x86\_64 ext4 combined image was downloaded from the official OpenWrt release repository. Since VMware does not natively boot raw IMG files efficiently, the OpenWrt disk image was converted into VMware VMDK format prior to deployment.

\subsection{Virtual Machine Configuration}

A custom virtual machine configuration was selected in VMware Workstation Pro to allow manual control over networking interfaces and storage configuration.

\begin{table}[H]
\centering
\caption{Initial OpenWrt VM Configuration}
\begin{tabular}{|l|l|}
\hline
Setting & Value \\ \hline
Guest OS & Other Linux 5.x 64-bit \\ \hline
vCPUs & 1 \\ \hline
Memory & 256 MB \\ \hline
Primary Disk & OpenWrt VMDK \\ \hline
Network Adapter 1 & NAT \\ \hline
Network Adapter 2 & Host-only \\ \hline
\end{tabular}
\end{table}

\subsection{Network Interface Design}

Two virtual network adapters were configured for each virtual router node:

\begin{itemize}
    \item NAT adapter for WAN connectivity and package installation
    \item Host-only adapter for internal mesh backbone communication
\end{itemize}

This architecture allowed the virtual environment to simulate isolated routed mesh connectivity while preserving Internet access for package installation and updates.

\subsection{Initial Boot Validation}

The OpenWrt virtual machine booted successfully after the VMDK disk image was attached to the VMware virtual SCSI controller.

During startup, OpenWrt initialized both virtual Ethernet interfaces successfully:

\begin{lstlisting}
eth0 NIC Link is Up
eth1 NIC Link is Up
\end{lstlisting}

The successful initialization confirmed that VMware virtual networking devices were correctly recognized by the OpenWrt kernel.

\begin{figure}[H]
\centering
\includegraphics[width=0.45\textwidth]{figures/openwrt-first-init.png}
\caption{Successful OpenWrt boot and interface initialization.}
\label{fig:bootlog}
\end{figure}

\section{Network Troubleshooting and Recovery}

\subsection{Initial Connectivity Failure}

Initial deployment attempts failed to provide Internet connectivity.

Observed symptoms included:

\begin{itemize}
    \item Missing default routes
    \item Failed DNS resolution
    \item Inability to ping external IP addresses
    \item DHCP inconsistencies
\end{itemize}

Example failure:

\begin{verbatim}
ping -c 4 8.8.8.8
Network unreachable
\end{verbatim}

Investigation revealed that VMware host-only networking was unintentionally acting as the WAN DHCP source. Multiple adapter configurations caused interface ambiguity and routing instability.

Observed VMware virtual networks included:

\begin{itemize}
    \item VMnet1: 192.168.61.0/24
    \item VMnet8: 192.168.111.0/24
\end{itemize}

\subsection{Controlled Simplification Strategy}

To isolate the networking issue, the second adapter was temporarily removed and the deployment simplified to a single NAT interface.

After simplification, the OpenWrt node successfully received:

\begin{itemize}
    \item DHCP assignment
    \item Default gateway
    \item External DNS access
\end{itemize}

Working configuration:

\begin{verbatim}
default via 192.168.111.2 dev eth0
\end{verbatim}

Successful validation tests included:

\begin{verbatim}
ping -c 4 8.8.8.8
ping -c 4 google.com
\end{verbatim}

All tests completed successfully with zero packet loss.

\subsection{Package Management Failure Analysis}

Package management operations initially failed during APK repository access.

Observed errors included:

\begin{verbatim}
wget exited with error 1
unexpected end of file
\end{verbatim}

Further investigation revealed that BusyBox wget lacked HTTPS support:

\begin{verbatim}
wget https://downloads.openwrt.org/
HTTPS support not compiled in.
\end{verbatim}

This explained repeated failures retrieving APK repository metadata over HTTPS.

The issue represented a significant operational weakness because administrators were effectively forced to downgrade repository transport from HTTPS to HTTP in order to restore package functionality.

Potential security implications included:

\begin{itemize}
    \item Man-in-the-middle attacks
    \item Package tampering
    \item Repository spoofing
    \item Supply-chain compromise
\end{itemize}

\subsection{Temporary Package Recovery}

HTTPS functionality was partially restored by manually transferring APK packages into the virtual machine and installing them locally:

\begin{verbatim}
apk add --allow-untrusted ./*.apk
\end{verbatim}

Repository transport was later temporarily modified from HTTPS to HTTP to restore repository index access.

After modification:

\begin{verbatim}
apk update
\end{verbatim}

successfully completed.

Core administrative packages were then installed successfully:

\begin{verbatim}
apk add nano
apk add vim
apk add curl
apk add tcpdump
apk add htop
apk add bash
apk add luci
\end{verbatim}

\section{LuCI Web Interface Deployment}

After installation of LuCI, the OpenWrt web interface was initially inaccessible through:

\begin{verbatim}
http://192.168.1.1
\end{verbatim}

Diagnostic analysis using:

\begin{verbatim}
netstat -tlnp
\end{verbatim}

confirmed that the web server was actively listening on ports 80 and 443:

\begin{verbatim}
0.0.0.0:80 LISTEN
0.0.0.0:443 LISTEN
\end{verbatim}

Further analysis revealed that the Windows host system could only directly access the VMware NAT interface rather than the internal OpenWrt LAN bridge.

The LuCI interface became accessible using:

\begin{verbatim}
http://192.168.111.130
\end{verbatim}

rather than the internal bridge address.

This demonstrated the distinction between:

\begin{itemize}
    \item Internal OpenWrt LAN addressing
    \item Externally reachable VMware NAT addressing
    \item Hypervisor network segmentation behavior
\end{itemize}

\begin{figure}[H]
\centering
\includegraphics[width=0.45\textwidth]{figures/luci1.png}
\caption{Accessible LuCI web management login page}
\label{fig:luci1}
\end{figure}
\begin{figure}[H]
\centering
\includegraphics[width=0.45\textwidth]{figures/luci2.png}
\caption{Accessible LuCI web management interface through VMware NAT networking.}
\label{fig:luci2}
\end{figure}

\section{Multi-Node Mesh Backbone Construction}

\subsection{Dual-Interface Mesh Backbone}

After establishing a stable baseline, a second virtual adapter was reintroduced using VMware host-only networking.

The second interface was assigned as a dedicated mesh backbone network:

\begin{verbatim}
uci set network.mesh=interface
uci set network.mesh.proto='static'
uci set network.mesh.device='eth1'
uci set network.mesh.ipaddr='10.10.10.1'
uci set network.mesh.netmask='255.255.255.0'
uci commit network
/etc/init.d/network restart
\end{verbatim}

Resulting topology:

\begin{table}[H]
\centering
\caption{Dual-Interface OpenWrt Topology}
\begin{tabular}{|l|l|l|}
\hline
Interface & Address & Purpose \\ \hline
eth0 & 192.168.111.x & VMware NAT / WAN \\ \hline
eth1 & 10.10.10.x & Internal mesh backbone \\ \hline
\end{tabular}
\end{table}

Routing verification confirmed correct segmentation:

\begin{verbatim}
default via 192.168.111.2 dev eth0
10.10.10.0/24 dev eth1
\end{verbatim}

\begin{figure}[H]
\centering
\includegraphics[width=0.45\textwidth]{figures/node1-route.png}
\caption{Validated OpenWrt Node 1 routing table showing WAN and mesh backbone separation.}
\label{fig:ip-route}
\end{figure}


\subsection{VM Cloning Strategy}

Instead of reinstalling OpenWrt repeatedly, a stable “golden image” virtual machine was cloned using VMware full-clone functionality.

Resulting nodes included:

\begin{itemize}
    \item OpenWrt-Node1
    \item OpenWrt-Node2
    \item OpenWrt-Node3
    \item OpenWrt-Node4
\end{itemize}

Each node received:

\begin{itemize}
    \item Unique hostname
    \item Unique mesh backbone address
    \item VMware-generated MAC address
\end{itemize}

\subsection{Four-Node Mesh Backbone Validation}

The resulting mesh backbone addressing scheme included:

\begin{table}[H]
\centering
\caption{Four-Node Mesh Backbone Addressing}
\begin{tabular}{|l|l|}
\hline
Node & Mesh Address \\ \hline
OpenWrt-Node1 & 10.10.10.1 \\ \hline
OpenWrt-Node2 & 10.10.10.2 \\ \hline
OpenWrt-Node3 & 10.10.10.3 \\ \hline
OpenWrt-Node4 & 10.10.10.4 \\ \hline
\end{tabular}
\end{table}

Connectivity testing was performed between all nodes using ICMP echo requests.

Successful bidirectional communication was verified between all nodes with:

\begin{itemize}
    \item Zero packet loss
    \item Stable latency
    \item Successful route propagation
\end{itemize}

The successful validation demonstrated:

\begin{itemize}
    \item Stable Layer 2 adjacency
    \item Functional Layer 3 routing
    \item Operational VMware host-only backbone networking
    \item Successful multi-node OpenWrt virtualization
\end{itemize}

\begin{figure}[H]
\centering
\includegraphics[width=0.45\textwidth]{figures/node4-ping.png}
\caption{Successful four-node mesh backbone connectivity validation.}
\label{fig:node4-ping}
\end{figure}

\subsection{Sequential Startup and DHCP Convergence}

Initial simultaneous startup of cloned nodes produced temporary WAN IP duplication behavior within the VMware NAT network.

Multiple nodes briefly inherited identical DHCP leases from previous cloned states.

This behavior was mitigated by sequential node startup. Sequential startup allowed VMware DHCP services to converge properly and assign unique WAN addresses based on MAC address differentiation.

Final WAN assignments stabilized as:

\begin{itemize}
    \item Node1: 192.168.111.134
    \item Node2: 192.168.111.136
    \item Node3: 192.168.111.137
    \item Node4: 192.168.111.138
\end{itemize}

This demonstrated:

\begin{itemize}
    \item DHCP lease convergence behavior
    \item Temporary ARP instability during cloned startup
    \item Virtualization timing dependencies
    \item Successful MAC-based lease differentiation
\end{itemize}

\subsection{IPv6 Duplicate Address Detection Behavior}

During sequential startup of cloned nodes, OpenWrt repeatedly reported IPv6 duplicate address detection warnings:

\begin{verbatim}
IPv6 duplicate address detected
\end{verbatim}

Despite repeated IPv6 duplicate detection events:

\begin{itemize}
    \item IPv4 routing remained stable
    \item Mesh backbone communication remained operational
    \item Node-to-node connectivity was unaffected
\end{itemize}

This behavior represents an important operational consideration for cloned embedded router deployments and virtualized distributed networking systems.

\begin{figure}[H]
\centering
\includegraphics[width=0.45\textwidth]{figures/ipv6error.png}
\caption{Observed IPv6 duplicate address detection during sequential startup of cloned OpenWrt nodes.}
\label{fig:ipv6error}
\end{figure}

\section{BATMAN-adv Mesh Routing Deployment}

\subsection{BATMAN-adv Installation}

After the successful construction of the four-node OpenWrt backbone topology, the BATMAN-adv routing protocol was deployed to provide dynamic Layer 2 mesh routing functionality.

The following packages were installed on all nodes:

\begin{verbatim}
apk add kmod-batman-adv batctl luci-proto-batman-adv
\end{verbatim}

Successful installation loaded the BATMAN kernel module dynamically:

\begin{verbatim}
lsmod | grep batman
\end{verbatim}

The BATMAN mesh interface was manually initialized:

\begin{verbatim}
modprobe batman-adv
ip link add name bat0 type batadv
ip link set up dev bat0
\end{verbatim}

The dedicated mesh transport interface was attached to BATMAN:

\begin{verbatim}
batctl meshif bat0 if add eth1
\end{verbatim}

\begin{figure}[H]
\centering
\includegraphics[width=0.48\textwidth]{figures/bat0add.png}
\caption{Successful initialization of the BATMAN-adv virtual mesh interface.}
\label{fig:bat0add}
\end{figure}

\subsection{Dynamic Neighbor Discovery}

BATMAN dynamically discovered neighboring routers using:

\begin{verbatim}
batctl meshif bat0 n
\end{verbatim}

Originator propagation tables were observed using:

\begin{verbatim}
batctl meshif bat0 o
\end{verbatim}

\begin{figure}[H]
\centering
\includegraphics[width=0.48\textwidth]{figures/mesh-routes.png}
\caption{BATMAN-adv originator routing tables dynamically populated across the mesh.}
\label{fig:meshroutes}
\end{figure}

\subsection{Four-Node Mesh Convergence}

All four OpenWrt nodes successfully joined the distributed BATMAN mesh.

The resulting deployment demonstrated:

\begin{itemize}
    \item Autonomous neighbor discovery
    \item Dynamic route propagation
    \item Persistent distributed convergence
    \item Self-healing mesh behavior
\end{itemize}

\begin{figure}[H]
\centering
\includegraphics[width=0.49\textwidth]{figures/pre-fail-top.png}
\caption{Stable four-node BATMAN-adv mesh topology before failure testing.}
\label{fig:prefail}
\end{figure}

\subsection{Node Failure and Reconvergence}

One node was intentionally powered off during operation to evaluate mesh resilience behavior.

Immediately after failure:

\begin{itemize}
    \item Neighbor tables updated dynamically
    \item Stale originators were removed
    \item Remaining nodes maintained connectivity
\end{itemize}

\begin{figure}[H]
\centering
\includegraphics[width=0.49\textwidth]{figures/post-fail-top.png}
\caption{BATMAN topology immediately following intentional node failure.}
\label{fig:postfail}
\end{figure}

\subsection{Node Reintegration}

After reboot, the failed node automatically rejoined the distributed mesh without manual reconfiguration.

\begin{figure}[H]
\centering
\includegraphics[width=0.49\textwidth]{figures/post-fail-top-post-boot.png}
\caption{Automatic BATMAN mesh reconvergence after node restoration.}
\label{fig:postboot}
\end{figure}

\subsection{Persistent BATMAN Configuration}

Persistent deployment was implemented using OpenWrt UCI configuration:

\begin{verbatim}
uci set network.bat0=interface
uci set network.bat0.proto='batadv'
uci set network.bat0.routing_algo='BATMAN_IV'

uci set network.mesh_bat=interface
uci set network.mesh_bat.proto='batadv_hardif'
uci set network.mesh_bat.master='bat0'
uci set network.mesh_bat.device='eth1'

uci commit network
\end{verbatim}

After reboot, all nodes automatically:
\begin{itemize}
    \item Recreated bat0
    \item Reattached eth1
    \item Rediscovered neighbors
    \item Rejoined the mesh
\end{itemize}

\subsection{BATMAN Native Ping and Traceroute}

BATMAN-native connectivity was validated using:

\begin{verbatim}
batctl meshif bat0 ping <neighbor MAC>
\end{verbatim}

Traceroute functionality was validated using:

\begin{verbatim}
batctl meshif bat0 traceroute <neighbor MAC>
\end{verbatim}

\begin{figure}[H]
\centering
\includegraphics[width=0.48\textwidth]{figures/traceroute.png}
\caption{BATMAN-native traceroute across the distributed mesh topology.}
\label{fig:traceroute}
\end{figure}

\subsection{MTU Optimization and Fragmentation Analysis}

BATMAN generated warnings indicating insufficient MTU allocation on the transport interface.

Initial configuration:

\begin{itemize}
    \item eth1 MTU: 1500
    \item bat0 MTU: 1500
\end{itemize}

The transport interface MTU was increased experimentally:

\begin{verbatim}
ip link set dev eth1 mtu 1532
\end{verbatim}

Following optimization:
\begin{itemize}
    \item Neighbor discovery remained stable
    \item BATMAN forwarding persisted
    \item Packet loss remained zero
\end{itemize}

Interestingly, bat0 continued operating at MTU 1500 while eth1 operated at MTU 1532.

\begin{figure}[H]
\centering
\includegraphics[width=0.48\textwidth]{figures/old-mtu.png}
\caption{Initial BATMAN deployment with insufficient MTU allocation.}
\label{fig:oldmtu}
\end{figure}

\begin{figure}[H]
\centering
\includegraphics[width=0.48\textwidth]{figures/new-mtu.png}
\caption{Stable BATMAN operation following MTU optimization.}
\label{fig:newmtu}
\end{figure}

\subsection{Management-Plane Addressing Issues}

During virtual-machine cloning, multiple nodes temporarily inherited overlapping management-plane identities.

Observed effects included:
\begin{itemize}
    \item Duplicate DHCP assignments
    \item Ambiguous LuCI access
    \item IPv6 duplicate-address warnings
\end{itemize}

Despite management-plane ambiguity, the BATMAN forwarding plane remained fully operational.

\begin{figure}[H]
\centering
\includegraphics[width=0.48\textwidth]{figures/nodeip-issue.png}
\caption{Management-plane addressing conflicts during cloned OpenWrt deployment.}
\label{fig:nodeipissue}
\end{figure}

\section{Configuration Export and Deployment Reproducibility}

To satisfy reproducibility requirements, configuration artifacts and operational mesh-state information were exported directly from all OpenWrt nodes.

The exported artifacts included:

\begin{itemize}
    \item Persistent OpenWrt UCI configuration files
    \item BATMAN neighbor tables
    \item BATMAN originator tables
    \item Routing tables
    \item Interface-state information
    \item Installed package inventories
\end{itemize}

A reusable export automation script was developed to simplify artifact extraction across all nodes.

The script automatically:
\begin{itemize}
    \item Collected configuration files
    \item Captured BATMAN operational state
    \item Archived deployment artifacts
    \item Generated compressed export bundles
\end{itemize}

The resulting deployment archives were extracted remotely from the OpenWrt nodes using SSH and SCP.

\subsection{Automated Configuration Export Script}

The following script was deployed on all OpenWrt nodes:

\begin{verbatim}
#!/bin/sh

NODE_NAME="$(cat /proc/sys/kernel/hostname)"
OUTDIR="/root/config-export"
ARCHIVE="/root/${NODE_NAME}-config-export.tar.gz"

mkdir -p "$OUTDIR"

cp /etc/config/network "$OUTDIR/network-${NODE_NAME}"
cp /etc/config/system "$OUTDIR/system-${NODE_NAME}"
cp /etc/config/firewall "$OUTDIR/firewall-${NODE_NAME}"
cp /etc/config/dhcp "$OUTDIR/dhcp-${NODE_NAME}"

batctl meshif bat0 if > "$OUTDIR/batman-if-${NODE_NAME}.txt"
batctl meshif bat0 n > "$OUTDIR/batman-neighbors-${NODE_NAME}.txt"
batctl meshif bat0 o > "$OUTDIR/batman-originators-${NODE_NAME}.txt"

ip addr > "$OUTDIR/ip-addr-${NODE_NAME}.txt"
ip route > "$OUTDIR/ip-route-${NODE_NAME}.txt"

apk list --installed > "$OUTDIR/packages-${NODE_NAME}.txt"

tar -czf "$ARCHIVE" -C /root config-export
\end{verbatim}

\subsection{Windows and Linux Script Compatibility Issues}

During deployment automation, shell scripts transferred from Windows systems to OpenWrt routers initially failed with:

\begin{verbatim}
-ash: /root/export-config.sh: not found
\end{verbatim}

Investigation revealed that the script contained Windows CRLF line endings rather than Linux LF line endings.

This caused BusyBox ash to misinterpret the script interpreter line:

\begin{verbatim}
#!/bin/sh^M
\end{verbatim}

The issue was corrected using:

\begin{verbatim}
sed -i 's/\r$//' /root/export-config.sh
\end{verbatim}

This represented an important interoperability consideration between Windows administrative systems and embedded Linux networking environments.

\subsection{SSH and SCP Artifact Extraction}

Configuration archives were remotely extracted from OpenWrt nodes using SCP over SSH.

Modern Windows SCP implementations attempted to default to SFTP mode, which was unsupported by the lightweight OpenWrt Dropbear deployment.

Observed error:

\begin{verbatim}
subsystem request failed on channel 0
scp: Connection closed
\end{verbatim}

Successful extraction required forcing legacy SCP protocol mode:

\begin{verbatim}
scp -O root@192.168.111.x:/root/archive.tar.gz .
\end{verbatim}

This demonstrated another operational limitation common within lightweight embedded Linux distributions where full SFTP subsystems are often omitted to conserve storage space and memory resources.

\subsection{Future Inter-Group Mesh Expansion}

Future work may interconnect independently deployed BATMAN mesh clusters between project groups using encrypted WireGuard tunnels.

Potential architecture:

\begin{itemize}
    \item Independent BATMAN mesh clusters
    \item WireGuard encrypted overlay tunnels
    \item Distributed inter-group route propagation
    \item Geographically separated mesh federation
\end{itemize}

This architecture would allow BATMAN-adv to operate across encrypted Internet-based transport overlays while preserving autonomous distributed mesh behavior.

\section{Future Raspberry Pi Deployment}

The final implementation stage will migrate the virtualized mesh environment onto Raspberry Pi hardware devices running OpenWrt and LibreMesh.

The Raspberry Pi deployment will evaluate:

\begin{itemize}
    \item Physical wireless mesh formation
    \item Real-world throughput and latency
    \item Hardware limitations
    \item Deployment obstacles
    \item Inter-group node interoperability
\end{itemize}

\section{Unresolved Operational Challenges and Future Investigation}

Although the distributed BATMAN-adv mesh deployment operated successfully overall, several operational inconsistencies and unresolved behaviors were identified during experimentation.

These issues represent valuable future investigation areas for both virtualized and physical mesh deployments.

\subsection{Persistent VMware DHCP Address Duplication}

Despite successful BATMAN mesh convergence, VMware NAT DHCP services occasionally assigned identical WAN-side addresses to multiple cloned OpenWrt nodes.

Example observed behavior included:

\begin{itemize}
    \item Multiple nodes simultaneously assigned 192.168.111.141
    \item Ambiguous SSH access
    \item Non-deterministic LuCI management behavior
    \item Administrative-plane collisions
\end{itemize}

Although sequential startup reduced the frequency of collisions, duplicate assignments still occasionally reappeared during extended experimentation.

The issue appeared isolated to VMware NAT-side DHCP behavior and did not affect BATMAN forwarding functionality operating on the dedicated mesh backbone interface.

This behavior represents an important virtualization-management limitation requiring further investigation.

\subsection{IPv6 Duplicate Address Detection Events}

OpenWrt repeatedly reported IPv6 duplicate address detection warnings during sequential startup and node cloning operations:

\begin{verbatim}
IPv6 duplicate address detected
\end{verbatim}

Despite repeated warnings:

\begin{itemize}
    \item IPv4 routing remained operational
    \item BATMAN mesh convergence persisted
    \item Node-to-node communication remained stable
\end{itemize}

The project did not fully investigate:

\begin{itemize}
    \item IPv6 address autoconfiguration behavior
    \item Duplicate IPv6 neighbor advertisements
    \item IPv6 interactions with BATMAN-adv
    \item Potential virtualization-layer MAC/address reuse
\end{itemize}

Future work should evaluate whether disabling IPv6 entirely or implementing static IPv6 addressing would improve deployment stability.

\subsection{Management-Plane and Forwarding-Plane Separation}

The project revealed significant architectural separation between:

\begin{itemize}
    \item Administrative management-plane behavior
    \item BATMAN distributed forwarding behavior
\end{itemize}

Even during severe management-plane inconsistencies:

\begin{itemize}
    \item BATMAN neighbor discovery remained stable
    \item Originator propagation continued
    \item Distributed forwarding functionality persisted
\end{itemize}

This highlighted the resilience of the distributed forwarding plane while simultaneously exposing administrative fragility caused by virtualization and DHCP interactions.

\subsection{Lightweight Embedded Linux Limitations}

Several OpenWrt operational limitations were identified during experimentation:

\begin{itemize}
    \item Missing HTTPS support in BusyBox wget
    \item Missing SFTP subsystem support in Dropbear
    \item Minimal shell-tool availability
    \item CRLF/LF shell-script compatibility issues
\end{itemize}

These limitations complicated:
\begin{itemize}
    \item Package installation
    \item Secure artifact transfer
    \item Deployment automation
    \item Remote administration
\end{itemize}

While expected for lightweight embedded environments, these limitations introduced significant operational overhead during deployment and experimentation.

\subsection{Incomplete Multi-Hop Routing Validation}

Although the BATMAN mesh operated successfully across all four nodes, the VMware host-only topology allowed all routers to remain directly adjacent at Layer 2.

As a result, the deployment did not fully validate:

\begin{itemize}
    \item True multi-hop relay behavior
    \item Intermediate-node forwarding dependency
    \item Mesh partition recovery
    \item Indirect route optimization
\end{itemize}

Future work should intentionally isolate subsets of nodes to force multi-hop forwarding paths and evaluate true relay-based mesh behavior.

\subsection{Future WireGuard-Based Inter-Mesh Federation}

Future experimentation may connect independent BATMAN mesh deployments between project groups using encrypted WireGuard overlay tunnels.

Potential future architecture:

\begin{itemize}
    \item Independent local BATMAN clusters
    \item WireGuard encrypted Internet overlays
    \item Distributed inter-group route propagation
    \item Geographically separated mesh federation
\end{itemize}

This would allow evaluation of:
\begin{itemize}
    \item Distributed mesh scalability
    \item Internet-based overlay transport
    \item Secure inter-cluster communication
    \item Hybrid mesh-overlay architectures
\end{itemize}

\section{Conclusion}

This project successfully implemented a persistent four-node distributed OpenWrt mesh-routing laboratory environment using VMware virtualization and BATMAN-adv.

The completed deployment demonstrated:

\begin{itemize}
    \item Successful OpenWrt virtualization
    \item Stable multi-node routed mesh backbone construction
    \item Autonomous BATMAN neighbor discovery
    \item Dynamic distributed originator propagation
    \item Persistent mesh deployment through UCI configuration
    \item Autonomous reconvergence after node failure
    \item BATMAN-native ping and traceroute functionality
    \item MTU optimization and fragmentation analysis
    \item Packet-level mesh observation using tcpdump
    \item Identification of virtualization and management-plane edge cases
\end{itemize}

The project evolved significantly beyond initial router virtualization objectives and ultimately produced a functional distributed autonomous mesh-routing research platform suitable for continued LibreMesh experimentation and future Raspberry Pi deployment.

The resulting environment now supports:

\begin{itemize}
    \item Reproducible distributed networking experimentation
    \item Dynamic routing analysis
    \item Mesh-failure resilience evaluation
    \item Multi-node convergence observation
    \item Future wireless mesh deployment research
\end{itemize}

The completed environment represents a robust foundation for continued experimentation with decentralized community-networking infrastructure using LibreMesh and BATMAN-adv.

\bibliographystyle{IEEEtran}

\begin{thebibliography}{99}

\bibitem{openwrt}
OpenWrt Project, ``OpenWrt Documentation,'' [Online]. Available: \url{https://openwrt.org}

\bibitem{libremesh}
LibreMesh Project, ``LibreMesh Documentation,'' [Online]. Available: \url{https://libremesh.org}

\bibitem{bmx7}
BMX7 Routing Protocol, ``BMX7 GitHub Repository,'' [Online]. Available: \url{https://github.com/bmx-routing/bmx7}

\bibitem{batmanadv}
BATMAN-adv Project, ``B.A.T.M.A.N. advanced Documentation,'' [Online]. Available: \url{https://www.open-mesh.org/projects/batman-adv/wiki}

\bibitem{batctl}
Open-Mesh, ``batctl User Guide,'' [Online]. Available: \url{https://www.open-mesh.org/projects/batman-adv/wiki/Batctl}

\bibitem{vmware}
VMware Inc., ``VMware Workstation Pro Documentation,'' [Online]. Available: \url{https://docs.vmware.com}

\end{thebibliography}

\bibliographystyle{IEEEtran}
\appendices

\section{Exported Configuration and Operational Artifacts}

The following exported artifacts were generated and archived from each OpenWrt node:

\begin{itemize}
    \item network-nodeX
    \item system-nodeX
    \item firewall-nodeX
    \item dhcp-nodeX
    \item batman-if-nodeX.txt
    \item batman-neighbors-nodeX.txt
    \item batman-originators-nodeX.txt
    \item ip-addr-nodeX.txt
    \item ip-route-nodeX.txt
    \item packages-nodeX.txt
\end{itemize}

These artifacts provide reproducible deployment state information and operational validation for the completed distributed BATMAN-adv mesh environment.

\end{document}